\chapter{S0：常温内参标定阶段}
\section{阶段定位与成果}
S0 的目标是建立并冻结常温 ACC/GYRO 内参及坐标映射关系；本章\textbf{不是要求重新进行已完成的常温标定}，而是为长期指南保留该阶段需完成的工艺检查、测试清单和可追溯输入。当前上传源码体现了部分实现，但不代表每项已有合格的真实台架数据。
\section{S0 测试与工艺检查项目总表}
\small
\begin{longtable}{C{17mm}p{48mm}p{50mm}p{28mm}}
\toprule 编号 & 检查／标定项目 & 执行要点与阶段产出 & 当前处理\\\midrule\endfirsthead
\toprule 编号 & 检查／标定项目 & 执行要点与阶段产出 & 当前处理\\\midrule\endhead
S0-01 & 转台与 IMU 配置确认 & 型号、坐标轴、量程、ODR、时间戳和工装记录 & 工艺前置\\
S0-02 & ACC 零偏与三轴全矩阵拟合 & 已知静止姿态、求解全矩阵及零偏 & 代码已有；实机核实\\
S0-03 & ACC 模型选择与留出验证 & 拟合/选模数据分离、条件数、残差 & 代码已有；留出验收需独立\\
S0-04 & GYRO 静止零速数据采集 & 开机/热稳态、多姿态零速基线 & 代码相关工艺已有\\
S0-05 & GYRO 正反向多速率采样 & 各轴已审核的低速工况、稳定窗口 & 代码已有；速率计量有限\\
S0-06 & GYRO 全矩阵与零偏拟合 & 方向覆盖、拟合条件数、候选模型比较 & 需核对软件输出和现场数据\\
S0-07 & 工艺数据质量检查 & 丢帧、姿态到位、稳态、异常值与重试审计 & 区分质量阈值/精度指标\\
S0-08 & 参数固化与版本追溯 & 矩阵定义、坐标系、CRC、参数/固件绑定 & S1 输入必须冻结\\
S0-09 & 标定结束前的初步合理性检查 & 模长、静止零速、方向符号、模型改善趋势 & 不能代替 S1 独立验收\\\bottomrule
\end{longtable}
\normalsize
\section{统一执行与交付要求}
\begin{enumerate}
\item 建立适用于当前转台安全运动范围的姿态与角速度工况，并在实际运动前确认外轴 X、内轴 Z 的坐标假设。
\item 保留原始数据、真实转台反馈、脚本版本、拟合残差、异常重试记录，明确\texttt{fit} 与\texttt{model\_select} 数据集身份。
\item 输出 $\mathbf C_a,\bm b_a,\mathbf C_g,\bm b_g$ 等实际采用的参数表示以及温度、量程、ODR、坐标系约定，完成不可变版本封存。
\item 完成上述 S0 后进入 S1，重新采集独立验收数据，不把参与拟合或模型选择的数据直接重用为“最终独立精度验收”。
\end{enumerate}
\begin{keynote}
\textbf{本章细化状态：} 目前是从上传的标定工程归纳出的工艺项目清单，后续若需要回溯标定算法，可再将 S0-01--S0-09 扩展为逐项详细作业指导。当前正文的成熟详细方案位于下一章 S1。
\end{keynote}
